La progettazione fisica ha l'obiettivo di tradurre lo schema logico relazionale in una struttura concretamente implementabile all'interno del DBMS PostgreSQL. Questa fase definisce rigorosamente i tipi di dato, le estensioni dei domini, i vincoli di integrità fisici e i meccanismi procedurali necessari a garantire il rispetto delle regole applicative precedentemente individuate.

Al fine di mantenere il capitolo compatto e focalizzato sulla discussione delle scelte ingegneristiche e architetturali, all'interno della trattazione sono stati isolati ed esaminati esclusivamente gli estratti di codice più significativi. La struttura completa della base di dati, comprensiva della creazione sequenziale di tutte le tabelle, delle chiavi esterne e del corpo integrale delle procedure memorizzate, è consultabile nella repository GitHub del progetto\footnote{\url{https://github.com/cossidente/database-lab-project/}}.

\section{Tipi enumerati e domini personalizzati}
\label{sec:tipi-enumerati-e-domini-personalizzati}

Per migliorare la qualità dei dati memorizzati e centralizzare i controlli di validazione, sono stati introdotti sia tipi enumerati sia domini personalizzati dotati di vincoli di controllo integrati.

I tipi enumerati sono stati impiegati per rappresentare insiemi finiti e statici di valori ammessi, impedendo l'inserimento di stringhe invalide e ottimizzando lo spazio di memorizzazione su disco, poiché PostgreSQL gestisce internamente tali strutture come identificatori numerici sequenziali.
In particolare, sono stati formalizzati i tipi \texttt{titolo\_enum}, \texttt{periodo\_enum}, \texttt{giorno\_enum} e \texttt{fascia\_oraria\_enum}.

Parallelamente, sono stati introdotti domini dedicati per vincolare gli attributi soggetti a restrizioni specifiche.
Il dominio \texttt{codice\_fiscale} verifica la correttezza formale del dato tramite un'espressione regolare strutturata per essere perfettamente corrispondente alla codifica standard ministeriale italiana di sedici caratteri alfanumerici, mentre \texttt{telefono\_dom} limita rigorosamente il formato dei recapiti a una stringa fissa di dieci cifre numeriche.

Un'attenzione particolare è stata riservata al dominio \texttt{anno\_accademico}. Oltre a validare la formattazione testuale nella configurazione \texttt{YYYY/YYYY}, la clausola di controllo estrae e converte le due componenti annuali per verificare che l'anno di chiusura sia esattamente successivo a quello di apertura, prevenendo l'inserimento di intervalli temporali incongruenti.
Analogamente, il dominio \texttt{punteggio} limita i voti d'esame all'intervallo discreto compreso tra 0 e 30.

L'implementazione fisica del dominio relativo all'anno accademico è stata realizzata attraverso la seguente struttura:

\begin{lstlisting}[language=SQL]
CREATE DOMAIN anno_accademico AS CHAR(9)
    CHECK (VALUE ~ '^[0-9]{4}/[0-9]{4}$' AND
           CAST(split_part(VALUE, '/', 2) AS INTEGER) =
           CAST(split_part(VALUE, '/', 1) AS INTEGER) + 1);
\end{lstlisting}

Questa scelta progettuale centralizza la logica di validazione nel server della base di dati, azzerando la necessità di replicare i controlli sul codice delle applicazioni esterne che interagiranno con l'infrastruttura dell'ateneo.

\subsection{Implementazione dei vincoli strutturali}
\label{ssec:implementazione-dei-vincoli-strutturali}

I vincoli strutturali derivati dalla progettazione logica sono stati mappati sfruttando le clausole native di chiave primaria, chiave esterna, vincoli di unicità e vincoli \texttt{CHECK} locali.

Le chiavi esterne garantiscono l'integrità referenziale del database: grazie alle regole di aggiornamento e cancellazione in cascata, impediscono la creazione di record orfani. I vincoli di unicità, d'altra parte, permettono di applicare logiche di business specifiche che la sola chiave primaria non basterebbe a gestire.

Un esempio cruciale è rappresentato dalla tabella \texttt{edizione}, all'interno della quale è stato introdotto un vincolo di unicità composto volto a escludere ridondanze annuali concorrenti:

\begin{lstlisting}[language=SQL]
CONSTRAINT uq_singola_edizione_annuale
    UNIQUE (codice_insegnamento, anno_accademico)
\end{lstlisting}

Questo costrutto garantisce in modo stringente il rispetto della specifica secondo cui per un dato insegnamento non può essere attivata più di una singola edizione all'interno dello stesso anno accademico.

L'introduzione di questo specifico vincolo si ricollega direttamente alle scelte di strutturazione discusse nella Sezione~\ref{sec:schema-logico}.
Come precedentemente illustrato, è proprio l'esplicita dichiarazione di questo vincolo di unicità annuale che permette di mantenere l'attributo \texttt{periodo} all'interno della chiave primaria della tabella \texttt{edizione}, al fine di favorirne la propagazione come chiave esterna verso la tabella \texttt{lezioni} e ottimizzare le successive operazioni di filtraggio.
La clausola \texttt{UNIQUE} agisce come un vincolo d'integrità a livello di DBMS, garantendo che la superchiave composta non minimale scelta per la chiave primaria non causi la violazione del vincolo di unicità annuale.

Ulteriori controlli di consistenza logica sono stati distribuiti tramite espressioni \texttt{CHECK} sulle singole tabelle. Tra i principali si segnalano il blocco preventivo sull'inserimento di valori negativi per i CFU nella tabella \texttt{insegnamento}, l'esclusione di relazioni autoreferenziali nella tabella \texttt{prerequisito} e la verifica che le date di verbalizzazione nella tabella \texttt{esame} non siano successive alla data odierna, prevenendo così l'inserimento di esiti d'esame con date future.

\subsection{Analisi delle scelte di indicizzazione e chiavi}
\label{ssec:analisi-delle-scelte-di-indicizzazione-e-chiavi}

Come anticipato nella progettazione logica, la scelta di includere l'attributo \texttt{periodo} all'interno della chiave primaria di \texttt{edizione}, nonostante il vincolo \texttt{UNIQUE(codice\_insegnamento, anno\_accademico)} sia già sufficiente a identificare univocamente ogni riga, risponde a una precisa strategia di ottimizzazione fisica.

Questa configurazione permette di propagare automaticamente il \texttt{periodo} come chiave esterna verso le tabelle dipendenti, come \texttt{lezione} e \texttt{insegnamento\_edizione}, portando un immediato vantaggio prestazionale.
Il DBMS ha sempre l'informazione sul periodo didattico localmente a disposizione all'interno di queste strutture, azzerando la necessità di eseguire continui e costosi \texttt{JOIN} con la tabella padre \texttt{edizione} ogni volta che si deve consultare il calendario settimanale o verificare la disponibilità di un'aula.

Il secondo grande beneficio riguarda la sicurezza dei dati, poiché la propagazione di una chiave composta crea un vincolo automatico a livello strutturale che impedisce a monte qualsiasi incongruenza temporale.
Se il \texttt{periodo} non facesse parte della chiave esterna, il database potrebbe accettare record formalmente validi ma logicamente contraddittori.

Per esempio, senza questo vincolo, una riga della tabella \texttt{lezione} potrebbe essere registrata per errore con il campo \texttt{periodo} impostato su \texttt{'2'} (secondo periodo), pur essendo collegata a un record della tabella \texttt{edizione} attivato esclusivamente nel \texttt{'1'} (primo periodo). Includendo il periodo all'interno della foreign key, il DBMS intercetta e rifiuta automaticamente inserimenti disallineati di questo tipo, garantendo l'integrità del calendario senza dover gravare sul sistema con trigger o controlli aggiuntivi.

\subsection{Ottimizzazione tramite indici}
\label{ssec:ottimizzazione-tramite-indici}

A completamento dell'analisi sulle chiavi, sono stati introdotti ulteriori indici (definiti nel file \texttt{sql/04\_indexes.sql}) per ottimizzare le operazioni più frequenti, distinguendo tre categorie di necessità.

\begin{lstlisting}[language=SQL]
-- Indici per i filtri di ricerca (Query)
CREATE INDEX idx_docente_cognome_nome
    ON docente (cognome, nome);
CREATE INDEX idx_insegnamento_nome
    ON insegnamento (nome);
CREATE INDEX idx_insegnamento_edizione_ricerca
    ON insegnamento_edizione (codice_insegnamento, anno_accademico, periodo);
CREATE INDEX idx_abilitazione_insegnamento
    ON abilitazione_docente_insegnamento (codice_insegnamento);

-- Indici per le seconde colonne di PK composite (Query + FK)
CREATE INDEX idx_piano_di_studio_insegnamento
    ON piano_di_studio (codice_insegnamento);
CREATE INDEX idx_esame_insegnamento
    ON esame (codice_insegnamento);

-- Navigazione inversa dei prerequisiti (CTE ricorsiva nel trigger)
CREATE INDEX idx_prerequisito_prerequisito
    ON prerequisito (codice_prerequisito);
\end{lstlisting}

La prima categoria riguarda i filtri di ricerca utilizzati dalle interrogazioni presentate nella Sezione~\ref{sec:operazioni}. In particolare, \texttt{idx\_do\-cen\-te\_co\-gno\-me\_no\-me} accelera la ricerca di un docente per nome e cognome (Sezione~\ref{ssec:estrazione-del-calendario-settimanale-di-un-docente}), \texttt{idx\_in\-seg\-na\-men\-to\_no\-me} velocizza il filtro per nome di un insegnamento (Sezione~\ref{ssec:media-dei-tentativi-per-gli-studenti-promossi}), mentre \texttt{idx\_in\-seg\-na\-men\-to\_e\-di\-zio\-ne\_ri\-cer\-ca} supporta la giunzione tra \texttt{le\-zio\-ne} ed \texttt{e\-di\-zio\-ne} tramite la tabella \texttt{in\-seg\-na\-men\-to\_e\-di\-zio\-ne}, sfruttata nella medesima operazione. A questi si aggiunge \texttt{idx\_a\-bi\-li\-ta\-zio\-ne\_in\-seg\-na\-men\-to} su \texttt{a\-bi\-li\-ta\-zio\-ne\_do\-cen\-te\_in\-seg\-na\-men\-to\-.co\-di\-ce\_in\-seg\-na\-men\-to}, utilizzato sia dal trigger \texttt{con\-trol\-la\_a\-bi\-li\-ta\-zio\-ne\_do\-cen\-te} sia dalla query di sostituzione cattedre (Sezione~\ref{ssec:sostituzione-in-blocco-di-una-cattedra-annuale}), che ricercano le abilitazioni per insegnamento anziché per docente.

La seconda categoria riguarda le foreign key che corrispondono alla seconda colonna di una chiave primaria composta, e che quindi non beneficiano dell'indice automaticamente creato sulla chiave primaria. È il caso di \texttt{e\-sa\-me\-.co\-di\-ce\_in\-seg\-na\-men\-to} e \texttt{pia\-no\_di\_stu\-dio\-.co\-di\-ce\_in\-seg\-na\-men\-to}, entrambe protette da un vincolo \texttt{ON~DELETE~RESTRICT} verso \texttt{in\-seg\-na\-men\-to}: senza un indice dedicato, ogni cancellazione di un insegnamento richiederebbe una scansione completa di queste tabelle (rispettivamente \num{320000} e \num{320000} tuple secondo la Tabella~\ref{tab:volumi}) per verificarne l'assenza di riferimenti. Gli indici \texttt{idx\_e\-sa\-me\_in\-seg\-na\-men\-to} e \texttt{idx\_pia\-no\_di\_stu\-dio\_in\-seg\-na\-men\-to} riducono tale controllo a un accesso diretto.

La terza categoria riguarda la navigazione inversa nella tabella \texttt{pre\-re\-qui\-si\-to}. Il trigger \texttt{con\-trol\-la\_ci\-clo\_pre\-re\-qui\-si\-ti} esegue una CTE ricorsiva che, a partire dal nuovo \texttt{co\-di\-ce\_pre\-re\-qui\-si\-to} inserito, risale la catena dei prerequisiti seguendo la colonna \texttt{co\-di\-ce\_pre\-re\-qui\-si\-to} stessa. Poiché la chiave primaria di \texttt{pre\-re\-qui\-si\-to} è definita su \texttt{(co\-di\-ce\_in\-seg\-na\-men\-to,\allowbreak\ co\-di\-ce\_pre\-re\-qui\-si\-to)}, le ricerche per \texttt{co\-di\-ce\_pre\-re\-qui\-si\-to} non sono coperte da alcun indice automatico: l'indice \texttt{idx\_pre\-re\-qui\-si\-to\_pre\-re\-qui\-si\-to} colma questa lacuna, rendendo ogni passo della ricorsione un accesso diretto anziché una scansione sequenziale.

\section{Trigger e vincoli procedurali}
\label{sec:trigger-e-vincoli-procedurali}

Le regole di business complesse che richiedono l'interrogazione coordinata di più tabelle non sono esprimibili tramite vincoli dichiarativi tradizionali. Per superare questi limiti strutturali, sono stati implementati vincoli procedurali attraverso trigger dedicati associati a funzioni scritte in linguaggio \texttt{PL/pgSQL}. Di seguito vengono discussi i trigger più significativi, direttamente derivanti dalle scelte progettuali illustrate nei capitoli precedenti; la lista completa è disponibile nella repository GitHub del progetto all'indirizzo \url{https://github.com/cossidente/database-lab-project/blob/main/sql/02_triggers.sql}.

\paragraph{Controllo delle propedeuticità d'esame}
Prima di consentire l'inserimento di una nuova riga nella tabella \texttt{esame}, il sistema deve verificare lo stato della carriera dello studente. La funzione \texttt{controlla\_prerequisiti} intercetta l'operazione in modalità \texttt{BEFORE INSERT} e verifica se la matricola ha superato tutte le materie definite come propedeutiche per l'insegnamento in esame. In caso di esito negativo, l'operazione viene interrotta sollevando un'eccezione che elenca esplicitamente i codici dei corsi mancanti.

\begin{lstlisting}[language=SQL]
CREATE OR REPLACE FUNCTION controlla_prerequisiti()
RETURNS TRIGGER AS $$
DECLARE
    mancanti TEXT;
BEGIN
    SELECT string_agg(p.codice_prerequisito, ', ')
    INTO mancanti
    FROM prerequisito p
    WHERE p.codice_insegnamento = NEW.codice_insegnamento
      AND NOT EXISTS (
          SELECT 1 FROM esame e
          WHERE e.matricola = NEW.matricola
            AND e.codice_insegnamento = p.codice_prerequisito
            AND e.punteggio >= 18
      );

    IF mancanti IS NOT NULL THEN
        RAISE EXCEPTION 'Prerequisiti mancanti: %', mancanti;
    END IF;

    RETURN NEW;
END;
$$ LANGUAGE plpgsql;
\end{lstlisting}

\paragraph{Prevenzione dei cicli nelle propedeuticità}
Per preservare la coerenza logica dell'offerta formativa, viene impedita la creazione di relazioni di propedeuticità circolari all'interno della tabella \texttt{prerequisito}. Il controllo viene eseguito mediante una query ricorsiva che esplora l'intero grafo diretto dei prerequisiti prima di autorizzare la transazione.

\begin{lstlisting}[language=SQL]
CREATE OR REPLACE FUNCTION controlla_ciclo_prerequisiti()
RETURNS TRIGGER AS $$
BEGIN
    IF EXISTS (
        WITH RECURSIVE raggiungibili(codice) AS (
            SELECT NEW.codice_prerequisito
            UNION
            SELECT p.codice_prerequisito
            FROM prerequisito p
            JOIN raggiungibili r ON r.codice = p.codice_insegnamento
        )
        SELECT 1 FROM raggiungibili WHERE codice = NEW.codice_insegnamento
    ) THEN
        RAISE EXCEPTION 'Inserimento creerebbe un ciclo nei prerequisiti tra % e %',
            NEW.codice_insegnamento, NEW.codice_prerequisito;
    END IF;
    RETURN NEW;
END;
$$ LANGUAGE plpgsql;
\end{lstlisting}

\paragraph{Automazione della ridondanza del saldo crediti}
Coerentemente con lo studio costi-benefici formalizzato in precedenza, l'attributo \texttt{crediti\_acquisiti} nella tabella \texttt{studente} viene gestito in modo interamente automatizzato. La correttezza di questo meccanismo dipende però da un invariante fondamentale: uno studente non deve poter registrare un secondo esame superato per lo stesso insegnamento, altrimenti i crediti verrebbero sommati più volte. Questo vincolo è garantito dal trigger \texttt{controlla\_esame\_gia\_superato}, che intercetta ogni \texttt{INSERT} sulla tabella \texttt{esame} e blocca l'operazione se esiste già un tentativo con punteggio $\ge 18$ per la medesima combinazione di matricola e insegnamento.

\begin{lstlisting}[language=SQL]
CREATE OR REPLACE FUNCTION controlla_esame_gia_superato()
RETURNS TRIGGER AS $$
BEGIN
    IF EXISTS (
        SELECT 1 FROM esame
        WHERE matricola = NEW.matricola
          AND codice_insegnamento = NEW.codice_insegnamento
          AND punteggio >= 18
    ) THEN
        RAISE EXCEPTION 'Lo studente % ha gia superato l''insegnamento %.',
            NEW.matricola, NEW.codice_insegnamento;
    END IF;
    RETURN NEW;
END;
$$ LANGUAGE plpgsql;
\end{lstlisting}

Solo una volta garantito questo invariante, la funzione \texttt{aggiorna\_crediti\_acquisiti} può incrementare il saldo in modo sicuro ad ogni inserimento con esito sufficiente ($\ge 18$).

\begin{lstlisting}[language=SQL]
CREATE OR REPLACE FUNCTION aggiorna_crediti_acquisiti()
RETURNS TRIGGER AS $$
BEGIN
    IF NEW.punteggio >= 18 THEN
        UPDATE studente
        SET crediti_acquisiti = crediti_acquisiti + (
            SELECT crediti FROM insegnamento WHERE codice = NEW.codice_insegnamento
        )
        WHERE matricola = NEW.matricola;
    END IF;
    RETURN NEW;
END;
$$ LANGUAGE plpgsql;
\end{lstlisting}

Specularmente, per garantire la consistenza del saldo anche a fronte di rimozioni di record, la funzione \texttt{aggiorna\_crediti\_post\_delete} agisce in modalità \texttt{AFTER DELETE} decrementando i crediti solo se l'esame rimosso presentava un esito sufficiente ($\ge 18$), grazie alla clausola \texttt{WHEN} applicata direttamente al trigger.

\begin{lstlisting}[language=SQL]
CREATE OR REPLACE FUNCTION aggiorna_crediti_post_delete()
RETURNS TRIGGER AS $$
BEGIN
    UPDATE studente
    SET crediti_acquisiti = crediti_acquisiti - OLD.crediti_esame
    WHERE matricola = OLD.matricola;
    RETURN OLD;
END;
$$ LANGUAGE plpgsql;
\end{lstlisting}